% Глава: 24-security
% Статус: задел (кратко), будет дополнено.
\chapter{Основы безопасности}
\label{ch:24-security}

\emph{Раздел будет дополнен в следующих версиях
пособия. Пока --- краткая вводка и ключевые темы.}

Безопасность здесь --- прежде всего безопасность доступа к
серверу лаборатории \fpath{pglab.nsu.ru} и к вашей ВМ на нём.

\section{Модель пользователя и привилегий}

Каждый процесс работает от имени пользователя
(\nameref{ch:05-permissions}); \cmd{sudo} даёт
временные права суперпользователя. Правила:
\begin{itemize}
  \item минимальные привилегии: запускаться под
    тем пользователем, который действительно
    нужен;
  \item не хранить секреты в коде и в
    \fpath{.gitignore}-не-защитённых местах;
  \item SSH-ключи предпочтительнее паролей
    (\nameref{ch:02-linux-access}).
\end{itemize}

\section{Что будет в полном разделе}
\begin{itemize}
  \item Файловые права и SELinux/AppArmor
    на практике;
  \item Защита SSH: отключение входа по паролю,
    AllowUsers, fail2ban;
  \item Работа с секретами (переменные окружения,
    файловые ключи, Vault-подход);
  \item Сетевая изоляция: firewall (nftables/
    ufw), группы безопасности;
  \item Обновления и патчи, управление зависимостями.
\end{itemize}

\section{Задачи для самопроверки}
\begin{enumerate}
  \item Почему SSH-ключи безопаснее паролей?
  \item Какие минимальные привилегии должны быть
    у сервиса в контейнере?
  \item Где \emph{не} следует хранить секреты?
\end{enumerate}
